%%%%%%%%%%%%%%%%%%%%%%%%%%%%%%%%%%%%%%%%%%%%%%%%%%%%%%%%%%%%%%%%%%%%%%%%%%%%%%%%%%%
%%%  Filename: main.tex
%%%  ---
%%%  Template for Master Thesis at FMIPA UGM
%%%  Created using ugmfmipa.cls
%%%  --- 
%%%  Original Written by DTETI UGM, updated to suit FMIPA's Template.
%%%  Only supports thesis (S2) proposal and thesis (S2) for now
%%%  
%%%  Developed by Khalilullah Al Faath, Jeremiah Manurung
%%%  Github repo: https://github.com/khalilullahalfaath/template-tesis-fmipa

%%%%%%%%%%%%%%%%%%%%%%%%%%%%%%%%%%%%%%%%%%%%%%%%%%%%%%%%%%%%%%%%%%%%%%%%%%%%%%%%%%%

%% Use option "bahasa" or "english" 
%%    to change the basic language used
%% User option "proposaltesis" or "tesis"
%% 	  to change the format
% \documentclass[<proposaltesis/tesis>,<bahasa/english>]{ugmfmipa}
\documentclass[proposaltesis,bahasa]{ugmfmipa}
%\documentclass[proposaltesis,bahasa]{ugmfmipa}
%======================================
% Information Input
%======================================
% Input author's name and ID number
\author{<<AUTHOR>>}{<<NIM>>}
% Input the thesis' title
\titleind{<<JUDUL THESIS>>}
\titleeng{<<THESIS TITLE>>}

% Program and the head of the program
\program{<<Nama Prodi>>}{<<Program coordinator>>}{<<NIP>>}
% Name of department head and NIP
\departmenthead{<<Head of the department>>}{<<NIP>>}
\major{<<Major>>}
\yearsubmit{<<TAHUN PENDADARAN>>}
\examdate{<<Exam date>>}
% Name of thesis supervisors/promotors
\addsupervisor{Dosen Pembimbing 1, S.Kom., M.Cs., PhD.}{<<NIP xxxxxx>>}
\addsupervisor{Dosen Pembimbing 2, S.Kom., M.Cs., PhD.}{<<NIP xxxxxx>>}
%\addsupervisor{<<Supervisor 3>>}{<<NIP>>}
% Name of examiners
%\addexaminer{<<Examiner 1>>}{<<NIP 1>>}
%\addexaminer{<<Examiner 2>>}{<<NIP 2>>}
%\addexaminer{<<Examiner 3>>}{<<NIP 3>>}
%\addexaminer{<<Examiner 4>>}{<<NIP 4>>}
%\addexaminer{<<Examiner 5>>}{<<NIP 5>>}
%\addexaminer{<<Examiner 6>>}{<<NIP 6>>}
%\addexaminer{<<Examiner 7>>}{<<NIP 7>>}
%\addexaminer{<<Examiner 8>>}{<<NIP 8>>}
%\addexaminer{<<Examiner 9>>}{<<NIP 9>>}

%======================================

% correct bad hyphenation here [example]
% \babelhyphenation[<<english/bahasa>>]{op-tical net-works semi-conduc-tor}
%% Uncomment block of code below to disable hyphenation
%\tolerance=1
%\emergencystretch=\maxdimen
%\hyphenpenalty=10000
%\hbadness=10000

\begin{document}
%======================================
% Create cover etc
%======================================

%---- COVER ----
%\printcover{sample/logougm.png}{Pendadaran/Tesis/Ringkasan Tesis*}
\printcover{sample/logougm.png}{Bachelor}
% *Choose one

%---- ENDORSEMENT PAGE ----
% Select endorsement page type. If you want to use your own PDF file,  
% 	use \printendorsementpdf, or if you want to use JPG file, use 
% 	\printendorsementjpg. Otherwise, use \printendorsement.
% 	Choose one only. Comment out unused command(s).
%
\cleardoublepage \phantomsection
\printendorsement
%\printendorsementpdf
%\printendorsementjpg{sample/scanned-endorsement.jpg}

%---- OPTIONAL: STATEMENT PAGE (uncomment if needed) ----
% \cleardoublepage \phantomsection
% \chapterstatement{contents/statement/statement}
%\chapterstatementjpg{sample/scanned-statement.jpg}

%---- OPTIONAL: DEDICATION PAGE (uncomment if needed) ----
% \cleardoublepage \phantomsection
% \chapterdedication{contents/dedication/dedication}

%---- OPTIONAL: PREFACE PAGE (uncomment if needed) ----
% \cleardoublepage \phantomsection
% \chapterpreface{contents/preface/preface}

%======================================
% Create Table of Contents, List of Figures, List of Tables
% <Do not change this part>
%======================================
\cleardoublepage \phantomsection
\thetoc
\onehalfspacing
\tableofcontents
\singlespacing
\cleardoublepage \phantomsection
\thelot
\listoftables
\cleardoublepage \phantomsection
\thelof
\listoffigures

%======================================

%---- NOMENCLATURE PAGE (optional - uncomment if needed) ----
% \cleardoublepage \phantomsection
% \chapternomenclature{contents/nomenclature/nomenclature}

%---- INTISARI PAGE----
\cleardoublepage \phantomsection
\chapterintisari{contents/abstract/intisari}

%---- ABSTRACT PAGE----
\cleardoublepage \phantomsection
\chapterabstract{contents/abstract/abstract}


%======================================



%======================================
%  MAIN TEXT
%======================================
\startmain

% --- Bab 1 & 2: Sama untuk proposal dan tesis ---
\inputChapter{contents/chapter-1/chapter-1}
\inputChapter{contents/chapter-2/chapter-2}

% --- Bab 3: Metodologi Penelitian (sama untuk proposal dan tesis) ---
\inputChapter{contents/chapter-3/chapter-3}

% --- Bab 4+: Berbeda untuk proposal dan tesis ---
\proposalContent{
    \inputChapter{contents/chapter-4/chapter-4-proposal}  % Bab 4: Jadwal Penelitian
}

\thesisContent{
    \inputChapter{contents/chapter-4/chapter-4}            % Bab 4: Implementasi
    \inputChapter{contents/chapter-5/chapter-5}            % Bab 5: Hasil Penelitian dan Pembahasan
    \inputChapter{contents/chapter-6/chapter-6}            % Bab 6: Kesimpulan dan Saran
}



%======================================

%======================================
%  References
%======================================
\cleardoublepage \phantomsection
\printbibliography[heading=bibintoc]

%======================================

%======================================
%  Appendix
%======================================
% You can change 
%    the filename and location of the files inputted
%    use \chapterappendix for the first page of the appendix
%    use \chapterappendixadd for the next page

\appendix
\appendixtables
\appendixfigures

\proposalContent{
    % Lampiran Proposal:
    % 1. Paper Referensi Utama
    \cleardoublepage \phantomsection
    \chapterappendix{contents/appendix/appendix-paper-referensi}
    % 2. Surat Pernyataan Penggunaan AI
    \chapterappendixadd{contents/appendix/appendix-surat-ai}
}

\thesisContent{
    % Lampiran Tesis:
    % 1. Surat Pernyataan Plagiasi
    \cleardoublepage \phantomsection
    \chapterappendix{contents/appendix/appendix-isi-lampiran}
    % 2. Surat Pernyataan Penggunaan AI
    \chapterappendixadd{contents/appendix/appendix-surat-ai}
    % 3. Bukti Submit Paper
    \chapterappendixadd{contents/appendix/appendix-bukti-submit}
    % 4. Paper Yang di Submit
    \chapterappendixadd{contents/appendix/appendix-paper-submit}
    % 5. Hasil Pengecekan Similaritas Paper
    \chapterappendixadd{contents/appendix/appendix-similaritas}
}




%======================================

\end{document}